\documentclass[14pt,a4paper]{extarticle}
\usepackage[a4paper,top=15mm,bottom=15mm,left=16mm,right=16mm]{geometry}
\usepackage[T1]{fontenc}
\usepackage{amsmath}
\usepackage{amssymb}
\usepackage{tgadventor}
\renewcommand*{\familydefault}{\sfdefault}
\usepackage[italic]{mathastext}
\usepackage{xcolor}
\usepackage{array}
\usepackage{tikz}
\usepackage[most]{tcolorbox}
\usetikzlibrary{arrows.meta}

\pagestyle{empty}
\setlength{\parindent}{0pt}
\setlength{\parskip}{6pt}
\AtBeginDocument{\raggedright}

\newtcolorbox{sol}[1]{colback=blue!2,colframe=blue!35,boxrule=0.6pt,
  colbacktitle=blue!8,coltitle=blue!50!black,
  left=6pt,right=6pt,top=4pt,bottom=4pt,
  title={\bfseries #1},
  toptitle=2pt,bottomtitle=2pt,
  before upper=\raggedright}

\begin{document}

\begin{center}
  {\LARGE\bfseries Always, Sometimes or Never?}\\[2pt]
  {\large Worked Solutions}
\end{center}

\vspace{4pt}

\begin{sol}{Q1. \(x^2\) is bigger than \(x\).}
Complete the table:
\[
\renewcommand{\arraystretch}{1.2}
\begin{array}{|c|c|c|}
\hline
x & x^2 & \text{Is } x^2 > x?\\
\hline
3 & 9 & \text{Yes}\\
\hline
0.5 & 0.25 & \text{No}\\
\hline
-2 & 4 & \text{Yes}\\
\hline
\end{array}
\]

When \(x=3\), \(x^2=9\), which is bigger than \(3\). So the statement is
\textbf{true} when \(x=3\).

When \(x=0.5\), \(x^2=0.25\), which is smaller than \(0.5\). So \(x=0.5\)
is a \textbf{counterexample}: the statement is \textbf{false} when
\(x=0.5\).

The statement is true for some values of \(x\) and false for others, so it
is \textbf{sometimes true}.
\end{sol}

\begin{sol}{Q2. \(10-x\) is smaller than \(10\).}
One possible choice of values is \(x=4\), \(x=0\) and \(x=-2\):
\[
\renewcommand{\arraystretch}{1.2}
\begin{array}{|c|c|c|}
\hline
x & 10-x & \text{Is } 10-x < 10?\\
\hline
4 & 6 & \text{Yes}\\
\hline
0 & 10 & \text{No}\\
\hline
-2 & 12 & \text{No}\\
\hline
\end{array}
\]

When \(x=4\), \(10-x=6\), so the statement is true.

When \(x=-2\), \(10-x=12\), so \(x=-2\) is a \textbf{counterexample}.

The statement is \textbf{sometimes true}: it is true for positive values
of \(x\), but false for \(x=0\) and for negative values of \(x\).
\end{sol}

\begin{sol}{Q3. Counterexamples to false statements}
\textbf{(a)} \textit{Every expression has a constant.}\\[2pt]
The expression \(4x+y\) has no constant term. So the statement is false.

\vspace{4pt}
\textbf{(b)} \textit{Every expression with \(3\) terms has \(3\) variables.}\\[2pt]
The expression \(2a+3b+1\) has \(3\) terms, but only \(2\) variables,
namely \(a\) and \(b\). So the statement is false.

\vspace{4pt}
\textbf{(c)} \textit{The term \(-3x\) always has a negative value.}\\[2pt]
When \(x=-2\), the term \(-3x\) has the value
\[
-3(-2)=6,
\]
which is not negative. So the statement is false.

\vspace{4pt}
\textbf{(d)} \textit{Every term has a coefficient bigger than \(1\).}\\[2pt]
The term \(x\) has coefficient \(1\), which is not bigger than \(1\).
So the statement is false. Another example is \(0.5a\), whose coefficient
is \(0.5\).
\end{sol}

\begin{sol}{Q4. Always, sometimes or never}
\textbf{(a)} \(x+x=2x\): \textbf{always true}.\\[2pt]
Reason: \(x+x\) means two lots of \(x\), which is \(2x\), for every value
of \(x\). For example, when \(x=3\), \(3+3=6=2\times 3\).

\vspace{6pt}
\textbf{(b)} \(5-x\) and \(x-5\) have the same value: \textbf{sometimes true}.\\[2pt]
True when \(x=5\), because \(5-5=0\) and \(5-5=0\).
False when \(x=0\), because \(5-0=5\) but \(0-5=-5\).

\vspace{6pt}
\textbf{(c)} \(x^2\) is negative: \textbf{never true}.\\[2pt]
Reason: \(x^2\) is a square, so \(x^2\ge 0\) for every real number \(x\).
For example, \(0^2=0\), which is not negative. No value of \(x\) makes
\(x^2\) negative.

\vspace{6pt}
\textbf{(d)} \(\dfrac{x}{2}\) is smaller than \(x\): \textbf{sometimes true}.\\[2pt]
True when \(x=4\), because \(\dfrac{4}{2}=2\), and \(2<4\).
False when \(x=-4\), because \(\dfrac{-4}{2}=-2\), and \(-2\) is not
smaller than \(-4\).
\end{sol}

\begin{sol}{Q5. Sam's ten examples}
Sam is wrong because testing ten values only shows that the statement is
true for those ten values. It does not show that the statement is true for
every value of \(x\). A single counterexample is enough to show that a
statement is not always true.

A counterexample is \(x=0\), because
\[
3x=3\times 0=0
\qquad\text{and}\qquad
x+1=0+1=1.
\]
Since \(0\) is not bigger than \(1\), the statement is false when \(x=0\).
So Sam's statement is not always true; it is sometimes true.
\end{sol}

\begin{sol}{Q6. A statement that is sometimes true}
One possible answer is:
\[
x+2>5.
\]
This is sometimes true.

It is true when \(x=4\), because
\[
4+2=6>5.
\]
A counterexample is \(x=1\), because
\[
1+2=3,
\]
and \(3\) is not bigger than \(5\). So the statement is sometimes true.
\end{sol}

\end{document}